\chapter{Vulnerabilidades de redes IoT IEEE 802.15.4} % Vulnerabilidades de redes IoT IEEE 802.15.4
\label{ataquesredesieee}
En este punto se resumen los ataques vigentes y aplicables contra redes basadas en el estándar IEEE 802.15.4, clasificándolos según las capas del modelos OSI.

\section{Ataques a la capa física} % Ataques a la capa física
Los ataques a la capa física cubren principalmente las interferencias de señal y el acceso físico a los dispositivos.
\subsection{Interferencia de señal y radio}
El objetivo de los ataques de interferencia de radio (radio jamming) o señal (signal jamming) es corromper las comunicaciones entre dispositivos dentro de un canal,  debido a las interferencias. La diferencia fundamental entre radio jamming y signal jamming, es que en el ataque radio jamming se realizan interferencias solo emitiendo señales, en lugar de enviar paquetes. En estos ataques básicamente se crean interferencias emitiendo señales a un canal, provocando una denegación de servicio.
\subsection{Manipulación de nodos}
Otro ataque a la capa física es la manipulación de nodos (nodes tampering). Este ataque se basa en el supuesto de que el atacante tiene acceso físico directo a los nodos, con lo que puede alterar su configuración. Teniendo esto en cuenta, el atacante puede acceder a la memoria de un nodo, capturar información privada (incluidos los datos criptográficos), comprometer la función del nodo o destruir totalmente el hardware. Las redes inalámbricas industriales de sensores IoT son vulnerables al acceso físico a los nodos, ya que generalmente estos se distribuyen en un área amplia.




\section{Ataques a la capa MAC} % Ataques a la capa MAC
Los ataques a la capa MAC se dirigen a las especificaciones de la capa de enlace de datos para lograr principalmente una denegación de servicio.
\subsection{Back-off manipulation}
La utilización de CSMA en el protocolo Zigbee, implica que se puede realizar una denegación de servicio, a los nodos legítimos de la red, si el canal se encuentra ocupado siempre que un nodo intenta transmitir la información. Si esto ocurre, el nodo espera un tiempo aleatorio que se amplía exponencialmente para un nodo que siempre encuentre el canal ocupado.
\subsection{Same nonce attack}
En la especificación 802.15.4, los dispositivos que operan en modo seguro y proporcionan servicios de control de acceso, mantienen una lista de control de acceso (ACL) que identifica los nodos para recibir datos. El formato de entrada de la ACL consiste en la dirección de destino, las opciones de modo seguro, la clave relacionada y el contenido nonce. Una lista de ACL de un remitente puede incluir dos entradas con la misma clave y el mismo nonce. Si el remitente utiliza las mismas claves y el mismo nonce dentro de dos transmisiones diferentes, un adversario que pueda obtener estos dos textos cifrados diferentes puede recuperar información útil sin conocer la clave.
\subsection{Ataques de repetición}
El mecanismo de protección contra repetición se realiza comprobando el contador de un mensaje reciente con el contador obtenido anteriormente. Si el contador reciente es igual o menor que el anterior, entonces la trama será rechazada. Sin embargo, debido a la poca capacidad de cómputo de los dispositivos, este mecanismo puede fallar cuando se retransmite de forma continuada el mismo paquete a la red.

\subsection{Envío Ack}
En medio de una transmisión entre dos nodos legítimos, un atacante puede escuchar los números de secuencia no encriptados de las tramas, corrompiendo las mismas al interferir en el momento de la recepción. Posteriormente, el atacante envía una trama falsa ACK cuyo número de secuencia está relacionado al nodo remitente.
\subsection{Ataque PanId}
Los nodos de una red Zigbee conocen el identificador del coordinador (PanId). Si existe más de un coordinador operando en el mismo POS (espacio operativo personal), podría ocurrir lo que se conoce como conflicto PANId. Si se produce dicho conflicto PanId, o bien el coordinador lo detecta a través de las balizas recibidas o es notificado por uno de los nodos miembros al recibir la señal de dos coordinadores con el mismo identificador.
Al ser notificado, el coordinador realiza el procedimiento de resolución de conflictos. Este mecanismo cubre principalmente los escaneos de canales y el procedimiento de realineación del coordinador que incluye elegir un nuevo PANId y transmitirlo a todos los nodos. Después de la resincronización, la red está lista para comunicarse de manera estable, lo que implica que la resolución del conflicto se ha completado.
\\Sin embargo, existe un escenario de ataque en el cual un usuario malicioso envía frecuentemente mensajes de notificación de conflicto falsos al coordinador, obligándolo a realizar el procedimiento de resolución de conflictos en bucle. Un atacante que genere mensajes de notificación de conflictos PanId configurando el campo relacionado en los marcos de mensajes, puede usar estos mensajes falsos para evitar o retrasar la comunicación entre los dispositivos y el coordinador.
\subsection{Ataque GTS}
De acuerdo con el estándar IEEE 802.15.4, GTS es la parte de la supertrama reservada para la transmisión de tramas desde dispositivos de baja latencia en el modo CSMA/CA ranurado. Un posible escenario de ataque ocurre cuando un atacante que ha logrado la sincronización con el coordinador al recibir mensajes de baliza, aprende los tiempos GTS del coordinador mediante la extracción del descriptor GTS, el cual se encuentra dentro del marco de la baliza e indica la longitud y el inicio del GTS en la supertrama. Después de obtener los tiempos GTS asignados, el adversario puede crear interferencias en cualquiera de esos momentos, causando la colisión y corrupción de los paquetes de datos emitidos entre el nodo legítimo y el coordinador.
\clearpage



\section{Ataques a la capa de red} % Ataques a la capa de red
Los ataques a la capa de red en IoT generalmente están diseñados para obstaculizar el mecanismo de selección de ruta o la estrategia de encaminamiento.
\subsection{Información falsa sobre una ruta}
Se ataca el proceso de descubrimiento de la ruta, anunciando que el camino más corto entre dos puntos es a través de un nodo malicioso. Proporcionando de esta forma datos de enrutamiento incorrectos a la red.
\subsection{Hello flood}
La transmisión de mensajes ''Hello'' (solicitudes de saludo) a cualquier nodo que se encuentre en la red, supone una denegación de servicio temporal en la red y afecta al proceso de selección de ruta.
\subsection{Ataque Sinkhole}
Un atacante puede generar y enviar un anuncio falso a los nodos vecinos, informando que hay una actualización en el encaminamiento de las tramas. De esta forma, los nodos envían los mensajes a través de él, capturando el tráfico de la red.
\subsection{Agujero de gusano}
Este ataque se lleva a cabo por al menos dos nodos ubicados estratégicamente en la red, los cuales hacen un túnel de baja latencia entre sí. Estos tiene la ruta más corta entre nodos, publicándolo al resto de nodos y capturando el tráfico de la red.
\subsection{Ataque Sybil}
Este ataque se basa en el envío de muchas identificaciones diferentes a la red por parte de un atacante, lo que aumenta su probabilidad de ser seleccionado en muchas rutas.




\clearpage
\subsection{Agujero negro}
Afecta a las rutas establecidas por la red. Un nodo malicioso usa su protocolo de enrutamiento para publicar que tiene la ruta más corta al nodo destino. El nodo malicioso siempre tiene la accesibilidad para responder a la solicitud de ruta, capturando los paquetes de los nodos y haciendo que todos o alguno de ellos se pierdan.



\section{Ataques a la capa de transporte} % Ataques a la capa de transporte
\subsection{Inundación TCP SYN}
Un atacante envía muchas solicitudes de conexión a un nodo legítimo, consumiéndose así los recursos del nodo (principalmente la memoria) para mantener esas conexiones innecesarias. Sin embargo, este ataque puede no ser efectivo si hay un mecanismo de defensa especificado en el protocolo.
\subsection{Ataque de desincronización}
Un atacante que escucha la conexión entre dos puntos finales, puede falsificar mensajes a cualquiera de ellos para que el nodo receptor solicite la retransmisión de mensajes. Si el atacante puede atacar los mensajes que llevan datos de control de conexión específicos en el momento adecuado, la sincronización entre los dos puntos finales podría perderse.



\section{Ataques a la capa de aplicación} % Ataques a la capa de aplicación
\subsection{Ataque Overwhelming}
Este ataque está relacionado con las aplicaciones de monitorización basadas en eventos, como la detección de movimiento, en las cuales los sensores activan una acción al detectar un evento. Uno o varios atacante pueden intentar saturar los nodos, lo que hará que la red reenvíe una gran cantidad de tráfico al receptor.
\subsection{Denegación de servicio basada en la ruta}
Este ataque alimenta algunos paquetes reproducidos en la red en los nodos hoja. Mediante el reenvío de estos paquetes al nodo receptor, se consumirán recursos de la red, principalmente el ancho de banda. El atacante también puede disminuir la vida útil de los dispositivos al hacer que estos consuman energía, mediante la retransmisión de paquetes irrelevantes.